\documentclass[11pt]{article}
\usepackage[margin=1in]{geometry}
\usepackage[T1]{fontenc}
\usepackage[utf8]{inputenc}
\usepackage{amsmath,amssymb,amsthm}
\usepackage[skip=\baselineskip,indent=0pt]{parskip}
\usepackage{enumitem}
\usepackage{hyperref}

\newtheorem*{theorem}{Theorem}

\newcommand{\PSh}{\mathsf{PSh}}
\newcommand{\EvSh}{\mathsf{Sh}^{\mathrm{ev}}}
\newcommand{\Ctx}{\mathsf{Ctx}}
\newcommand{\Cart}{\mathsf{C}}
\newcommand{\Vcat}{\mathsf{V}}
\newcommand{\Set}{\mathbf{Set}}
\newcommand{\Bool}{\mathsf{Bool}}
\newcommand{\Chu}{\operatorname{Chu}}
\newcommand{\op}{\mathrm{op}}
\newcommand{\RepoURL}{https://codeberg.org/doc2cog/hoapa/src/branch/main}
\newcommand{\RepoDir}[1]{\href{\RepoURL/#1}{\nolinkurl{#1/}}}
\newcommand{\RepoFile}[1]{\href{\RepoURL/#1}{\nolinkurl{#1}}}

\title{Event-Local Sheafification and a Linear/Non-Linear Adjunction\\
with Chu Coefficients}
\author{Hoapa project}
\date{June 2026 \\[2pt]
\normalsize Extended abstract --- machine-checked in cubical Agda}

\begin{document}
\maketitle

\begin{abstract}
We construct an event-local sheafification functor and use it to lift a Benton
linear/non-linear adjunction over a small site of finite causal interaction
contexts. The objects of the site are finite local views of a concurrent
interaction: unityped actors, the events in which they participate, and the
causal order on those events. The topology is generated by
decomposing each context into its single-event subcontexts. For any coefficient
category with finite products,
sheafification sends a presheaf to the product of its values on those
subcontexts and is left adjoint to the inclusion of event-local sheaves. Given a
Benton adjunction $F\dashv G$ between a cartesian coefficient category and a
$*$-autonomous coefficient category, this construction yields an adjunction
between cartesian-valued and linear-valued event-local sheaves. The cartesian
side is the freely copyable-and-discardable regime, the $*$-autonomous side is
the resource-sensitive regime with linear duals, and the induced comonad is the
sheaf-level linear exponential ``$!$''. The hypotheses are also
jointly realized in a concrete coefficient model: $\Vcat=\Chu(\Set,\Bool)$ has
the required finite products and a minimal $*$-autonomous interface, and the
Benton adjunction $\Set\rightleftarrows\Chu(\Set,\Bool)$ is constructed in Agda.
Separately, the postulate-free session syntax is interpreted into this Chu model
by an underlying functor. Thus the sheaf-level theorem is generic, but not
vacuous.
\end{abstract}

This note isolates one theorem from a larger programme that models concurrent
interaction systems, and the typed data produced by them, in the language of
sheaves and adjoint functors. The theorem starts from a deliberately austere
view of agency. Actors are unityped: mathematically, they are elements of a bare
set, and the formal system exposes no role, origin, or hidden state on which a
construction can branch. The structure lies instead in events, causal
dependence, incidence between events and actors, and the typed values assigned
to finite contexts of interaction. The result below says that a coherent global
assignment over such a context is freely reconstructed from its event-local
assignments. The system's global point of view is not imposed from outside; it
is the coherent reconstruction of its event-local standpoints.

The theorem has two parts. First, event-local observations freely generate a
coherent global object. Second, that reconstruction respects the
cartesian-linear structure of the values assigned to contexts: reusable data
live on the cartesian side, resource-sensitive session-like data live on the
linear side, and the Benton adjunction relates the two at the sheaf level.

\section*{The event-local site}

The actor model of concurrency~\cite{Hewitt1973} describes computation as a
collection of agents that interact solely by exchanging messages. We use only
the unityped aspect of that tradition: actors are elements of a bare set, and
the kernel exposes no role, origin, or tag on which a construction can branch.
Lamport's causal order~\cite{Lamport1978} supplies the temporal structure:
events are partially ordered by possible influence. Following
presheaf models of concurrency~\cite{CattaniWinskel1996}, these finite
observations are organized as a small site.

Fix a set $A$ of actors. A \emph{finite causal context} over $A$ consists of a
finite set of events, a partial order on those events, and an incidence relation
recording which actors participate in which events. The context is
\emph{supported}: every actor present in the context participates in at least one
event. A morphism of contexts is a causality-reflecting embedding, preserving
actors, events, and incidence, and both preserving and reflecting the causal
order. Write $\Ctx_A$ for the resulting category; we call it the actor-causal
site.

The topology used here is event-local. For an event $e$ of a context $F$, write
$F_e$ for the single-event subcontext at $e$, whose actors are exactly those
incident to $e$. The generating cover of $F$ is
\[
  \{\,F_e\longrightarrow F\,\}_{e\in E(F)}.
\]
No other covers are made primitive in this theorem. In particular, causal
refinement, higher overlaps, and richer compatibility conditions are deferred to
later topologies. The Agda development inherited the internal word ``point'' for
these single-event subcontexts; in this abstract we use \emph{event-local} for
the site-side notion and reserve ``points'' and ``states'' for Chu spaces below.

\section*{Event-local sheafification}

Let $\mathcal C$ be a coefficient category with finite products. A
$\mathcal C$-valued presheaf on the site is a functor
\[
  P : \Ctx_A^{\op}\longrightarrow\mathcal C,
\]
and we write $\PSh_A(\mathcal C)=[\Ctx_A^{\op},\mathcal C]$. An event-local
sheaf is a presheaf satisfying descent for the single-event covers above. The
formalization proves the expected product-locality characterization: a presheaf
is event-local precisely when its value on a context is determined by its values
on the single-event subcontexts of that context.

The inclusion
$\iota_{\mathcal C}:\EvSh_A(\mathcal C)\hookrightarrow\PSh_A(\mathcal C)$ has an
explicit left adjoint. For a presheaf $P$, define
\[
  a_{\mathcal C}(P)(F)=\prod_{e\in E(F)} P(F_e),
\]
with restrictions induced by the corresponding maps between single-event
subcontexts. This is sheafification in the restricted sense used here: the
reflector onto presheaves satisfying descent for the explicitly presented
event-local covers. It is not claimed to be the associated-sheaf functor for an
independently formalized Grothendieck coverage. The machine-checked result is
the adjunction
\[
  a_{\mathcal C}\dashv\iota_{\mathcal C}.
\]

\section*{The linear/non-linear lift}

Let $\Cart$ be a cartesian category and let $\Vcat$ be a $*$-autonomous
category~\cite{Barr1979}, both with finite products. A linear/non-linear
adjunction in the sense of Benton~\cite{Benton1995} is a monoidal adjunction
$F\dashv G:\Cart\rightleftarrows\Vcat$ from the cartesian world to the linear
one. The induced comonad on $\Vcat$ is the linear exponential ``$!$''. The Agda
interface used by the theorem is intentionally minimal: it records the two
functors, the adjunction, and the strong-monoidal comparison maps of $F$. The
proof consumes the adjunction and product structure; the omitted global
coherence and tensor-hom universal property are not unfolded.

Postcomposition lifts $F$ and $G$ to presheaves,
\[
  F_* : \PSh_A(\Cart)\to\PSh_A(\Vcat),
  \qquad
  G_* : \PSh_A(\Vcat)\to\PSh_A(\Cart),
\]
with $F_*\dashv G_*$ computed pointwise. The right adjoint preserves
event-local sheaves: descent for the single-event cover is transported across
the hom-set bijection $\Cart(X,GY)\cong\Vcat(FX,Y)$. The left adjoint is made to
land in sheaves by composing with the reflector. With
\[
  L=a_{\Vcat}\circ F_*\circ\iota_{\Cart},
  \qquad
  R=G_*|_{\EvSh_A(\Vcat)},
\]
we obtain the theorem.

\begin{theorem}
Let $\Cart$ be a cartesian category and $\Vcat$ a $*$-autonomous category, both
with finite products, and let $F\dashv G:\Cart\rightleftarrows\Vcat$ be a
Benton linear/non-linear adjunction. Then there is an adjunction
\[
  L \dashv R : \EvSh_A(\Cart)\rightleftarrows\EvSh_A(\Vcat)
\]
between cartesian-valued and linear-valued event-local sheaves on the
actor-causal site. The induced comonad $L\circ R$ is the sheaf-level linear
exponential.
\end{theorem}

\section*{A concrete Chu coefficient model}

The theorem above is generic in its coefficient categories. The new Chu
construction changes its status: the hypotheses are not merely assumed, but are
simultaneously inhabited in a concrete machine-checked model. For a dualizing
set $k$, a Chu object over $\Set$ is a triple
$(A^+,A^-,\langle-,-\rangle:A^+\times A^-\to k)$~\cite{Barr1996,Pratt1995}.
In the standard Chu terminology, $A^+$ is the carrier of points and $A^-$ the
carrier of states; this is distinct from the event-local subcontexts of the
site.

At $k=\Bool$, the Agda development constructs the category
$\Chu(\Set,\Bool)$, its finite products, and the formalization's minimal
$*$-autonomous interface. It also constructs the Benton adjunction
\[
  \Set\rightleftarrows\Chu(\Set,\Bool),
\]
with right adjoint $G(A)=A^+$ and left adjoint
$F(X)=(X,\Bool^X,\operatorname{ev})$. Supplying these data to the generic
event-local theorem yields a concrete witness of the lifted adjunction
$L\dashv R$. In this construction, the Benton linear/non-linear adjunction is
the cartesian-linear bridge: it relates reusable cartesian values to
resource-sensitive linear values. The Chu model supplies concrete coefficients
in which that bridge exists, so the sheaf-level adjunction is an instantiated
theorem rather than a purely conditional one.

\section*{Session syntax in the model}

A separate result connects the session syntax to this coefficient model. The
postulate-free syntax of session types and their structural morphisms, in the
linear-propositions tradition for session types~\cite{CairesPfenning2010,Wadler2014},
is interpreted into any target with the formalization's minimal $*$-autonomous
and finite-product interface by an object map $J0$, a duality theorem
$J0(S^{\perp_s})\cong (J0(S))^\perp$, and a morphism map $J1$. Instantiating
that interpretation at $\Chu(\Set,\Bool)$ gives an underlying functor in the
Cubical Agda category library from the session category to the Chu coefficient
category. This is intentionally weaker than full session semantics: it is not an
operational send/receive model, not a coherent monoidal or $*$-autonomous
session functor, and not a construction of the relevant cyclic operadic action
in Chu. The remaining action postulate in \texttt{Sessions.agda} is therefore
not discharged by this result.

\section*{Scope and significance}

Three classical lines meet in the theorem. It is a descent statement over a
small site, in the tradition of Grothendieck~\cite{SGA4,MacLaneMoerdijk1992}.
It is an instance of Lawvere's principle that adjointness organizes
foundational structure~\cite{Lawvere1969}: the global object is obtained by a
left adjoint, and its universal property is the freeness of the reconstruction.
And it lifts Benton's linear/non-linear adjunction from coefficient categories
to sheaves. The Chu instance adds non-vacuity: the general theorem is still
stated parametrically, but there is now a concrete postulate-free coefficient
model in which its hypotheses hold.

In the wider programme, this fixes a bounded piece of the operational side of a
larger conjectural architecture. Event-local observations freely generate a
coherent global object, that reconstruction is compatible with
resource-sensitive values, and one concrete Chu coefficient model realizes the
required value-side structure. Claims about richer causal descent, operational
session semantics, formal introduction of external data into the system, and
cross-site modal adjunctions belong to later stages.

\section*{Related work}

The theorem has siblings by facet rather than a single direct predecessor. The
local-to-global pattern is close to Abramsky and Brandenburger's
sheaf-theoretic account of contextuality~\cite{AbramskyBrandenburger2011},
where contextuality is detected as an obstruction to global sections, and to
later cohomological developments of that line~\cite{AbramskyEtAl2015}.
Event-based sheaf models of systems appear in work of Zardini, Spivak, Censi,
and Frazzoli~\cite{ZardiniEtAl2021}. Local views of distributed computation are
organized topologically in work of Goubault, Ledent, and
Rajsbaum~\cite{GoubaultLedentRajsbaum2018}.

On the linear-concurrency side, thin concurrent games and generalized species
provide a recent adjacent model~\cite{ClairambaultOlimpieriPaquet2023}, while
work on session types derived from linear logic clarifies the session-theoretic
side~\cite{Wadler2014,VanDenHeuvelPerez2024}. Licata, Shulman, and Riley's
fibrational framework for substructural and modal logics~\cite{LSR2017}, and
the modal homotopy type theory of Riley, Finster, and
Licata~\cite{RileyFinsterLicata2021}, are further siblings on the
type-theoretic side: they organize movement between logical regimes rather than
collapsing those regimes into one. These works share pieces of the structure
below, but do not construct the actor-causal site, event-local sheafification,
Benton lift, and Chu coefficient instance proved here.

\section*{Navigating the Agda formalization}

The repository is:
\begin{center}
\url{https://codeberg.org/doc2cog/hoapa}
\end{center}
The theorem is checked in the subtree:
\begin{center}
\RepoDir{formal/Concurrency}
\end{center}
The most direct way into the proof is to start from the theorem module:
\begin{center}
\RepoFile{formal/Concurrency/Hybrid/PointLinearNonlinearAdjunction.agda}
\end{center}
then follow the imports backwards through the reflector, the sheaf adapter, the
site geometry, and finally the coefficient models.

The formalization splits into three compact layers. The role of Agda here is
not only to record the statement: it checks that the reflector, lifted
adjunction, finite-product witnesses, and Chu instance inhabit the same
interfaces simultaneously.

\paragraph{Event-local reflector and lifted adjunction.}
Folder: \RepoDir{formal/Concurrency/Hybrid}
\begin{itemize}[leftmargin=1.4em,itemsep=2pt]
\item \RepoFile{formal/Concurrency/Hybrid/SupportedChangeOfBase.agda}
\item \RepoFile{formal/Concurrency/Hybrid/SupportedLinearNonlinearSheafAdapter.agda}
\item \RepoFile{formal/Concurrency/Hybrid/PointSheafAdjunction.agda}
\item \RepoFile{formal/Concurrency/Hybrid/PointLinearNonlinearAdjunction.agda}
\end{itemize}

\paragraph{Event-local site geometry.}
Folder: \RepoDir{formal/Concurrency/Sites}
\begin{itemize}[leftmargin=1.4em,itemsep=2pt]
\item \RepoFile{formal/Concurrency/Sites/EventPoint.agda}
\item \RepoFile{formal/Concurrency/Sites/EventSupported.agda}
\item \RepoFile{formal/Concurrency/Sites/PointCover.agda}
\item \RepoFile{formal/Concurrency/Sites/PointOverlap.agda}
\item \RepoFile{formal/Concurrency/Sites/SupportedFiniteFrameCategory.agda}
\end{itemize}

\paragraph{Concrete Chu coefficient stack.}
Folders:
\begin{itemize}[leftmargin=1.4em,itemsep=1pt]
\item \RepoDir{formal/Concurrency/Algebras}
\item \RepoDir{formal/Concurrency/Hybrid}
\end{itemize}
\begin{itemize}[leftmargin=1.4em,itemsep=2pt]
\item \RepoFile{formal/Concurrency/Algebras/Chu.agda}
\item \RepoFile{formal/Concurrency/Algebras/ChuFiniteProducts.agda}
\item \RepoFile{formal/Concurrency/Algebras/ChuMonoidal.agda}
\item \RepoFile{formal/Concurrency/Hybrid/ChuLinearNonlinear.agda}
\item \RepoFile{formal/Concurrency/Hybrid/ChuPointVProd.agda}
\item \RepoFile{formal/Concurrency/Hybrid/ChuPointExample.agda}
\end{itemize}
The final Agda witness is named \texttt{Chu-nonvacuity}; its fully expanded type
is omitted because writing it out makes the type-checker expand definitions
prohibitively.

\paragraph{Session syntax and Chu interpretation.}
Folder: \RepoDir{formal/Concurrency/Algebras}
\begin{itemize}[leftmargin=1.4em,itemsep=2pt]
\item \RepoFile{formal/Concurrency/Algebras/SessionSyntax.agda}
\item \RepoFile{formal/Concurrency/Algebras/SessionMorphismSyntax.agda}
\item \RepoFile{formal/Concurrency/Algebras/SessionObjectModel.agda}
\item \RepoFile{formal/Concurrency/Algebras/SessionChuModel.agda}
\end{itemize}

\section*{Limitations}

We are explicit about what is not done here.
\begin{itemize}[leftmargin=1.4em,itemsep=2pt]
\item The value-level adjunction $F\dashv G$ is assumed in the general theorem;
it is discharged only in the concrete Chu instance
$(\Set,\Chu(\Set,\Bool))$.
\item The Chu model uses the formalization's minimal $*$-autonomous interface.
The global monoidal coherence and tensor-hom universal property of Chu
categories are not fully formalized here.
\item The topology only imposes descent along single-event covers. Causal order,
overlaps of events, and richer compatibility conditions do not yet carry descent
content.
\item Event-local sheaves are defined by descent for the explicitly presented
single-event covers. We do not yet compare this with a fully formalized
Grothendieck coverage or library sheaf condition for the same site.
\item The Sessions-to-Chu interpretation is an underlying functorial
realization of the syntax, not operational send/receive semantics and not a
coherent monoidal or $*$-autonomous session semantics.
\item This is not a cross-site modal adjunction, not a construction of the
required cyclic operadic action in Chu, and not an account of the formal
introduction of external data into the system.
\end{itemize}

\section*{A note on authorship}

This abstract was drafted with the assistance of a language model and reviewed
by the authors; the surrounding framework and its philosophical motivation are
presented separately in the project repository at
\url{https://codeberg.org/doc2cog/hoapa}. The formal claims reported here are
backed by proofs that the cubical Agda type-checker accepts in the modules named
above, with the scope limitations stated explicitly.

\begin{thebibliography}{20}
\raggedright

\bibitem{Hewitt1973}
C. Hewitt, P. Bishop, and R. Steiger.
A universal modular ACTOR formalism for artificial intelligence.
In \emph{Proceedings of the 3rd International Joint Conference on Artificial
Intelligence (IJCAI)}, pp. 235--245, 1973.

\bibitem{Lamport1978}
L. Lamport.
Time, clocks, and the ordering of events in a distributed system.
\emph{Communications of the ACM} 21(7):558--565, 1978.

\bibitem{CattaniWinskel1996}
G. L. Cattani and G. Winskel.
Presheaf models for concurrency.
BRICS Report Series RS-96-35, 1996. Also in
\emph{Computer Science Logic (CSL '96)}, LNCS 1258, Springer, 1997.

\bibitem{Barr1979}
M. Barr.
\emph{$*$-Autonomous Categories}.
Lecture Notes in Mathematics 752. Springer, 1979.

\bibitem{Benton1995}
P. N. Benton.
A mixed linear and non-linear logic: proofs, terms and models.
In \emph{Computer Science Logic (CSL '94)}, LNCS 933, pp. 121--135. Springer,
1995.

\bibitem{Barr1996}
M. Barr.
The Chu construction.
\emph{Theory and Applications of Categories} 2(2):17--35, 1996.

\bibitem{Pratt1995}
V. R. Pratt.
The Stone gamut: a coordinatization of mathematics.
In \emph{Proceedings of the 10th Annual IEEE Symposium on Logic in Computer
Science (LICS)}, pp. 444--454. IEEE, 1995.

\bibitem{CairesPfenning2010}
L. Caires and F. Pfenning.
Session types as intuitionistic linear propositions.
In \emph{CONCUR 2010}, LNCS 6269, pp. 222--236. Springer, 2010.

\bibitem{Wadler2014}
P. Wadler.
Propositions as sessions.
\emph{Journal of Functional Programming} 24(2--3):384--418, 2014.

\bibitem{SGA4}
M. Artin, A. Grothendieck, and J.-L. Verdier, eds.
\emph{Th\'eorie des topos et cohomologie \'etale des sch\'emas (SGA 4)}.
Lecture Notes in Mathematics 269, 270, 305. Springer, 1972--1973.

\bibitem{MacLaneMoerdijk1992}
S. Mac Lane and I. Moerdijk.
\emph{Sheaves in Geometry and Logic: A First Introduction to Topos Theory}.
Springer, 1992.

\bibitem{Lawvere1969}
F. W. Lawvere.
Adjointness in foundations.
\emph{Dialectica} 23(3--4):281--296, 1969.

\bibitem{AbramskyBrandenburger2011}
S. Abramsky and A. Brandenburger.
The sheaf-theoretic structure of non-locality and contextuality.
\emph{New Journal of Physics} 13:113036, 2011.

\bibitem{AbramskyEtAl2015}
S. Abramsky, R. S. Barbosa, K. Kishida, R. Lal, and S. Mansfield.
Contextuality, cohomology and paradox.
In \emph{24th EACSL Annual Conference on Computer Science Logic (CSL 2015)},
LIPIcs 41, pp. 211--228. Schloss Dagstuhl, 2015.

\bibitem{ZardiniEtAl2021}
G. Zardini, D. I. Spivak, A. Censi, and E. Frazzoli.
A compositional sheaf-theoretic framework for event-based systems.
arXiv:2101.10485, 2021.

\bibitem{GoubaultLedentRajsbaum2018}
E. Goubault, J. Ledent, and S. Rajsbaum.
A simplicial complex model for dynamic epistemic logic to study distributed task
computability.
arXiv:1809.03095, 2018.

\bibitem{ClairambaultOlimpieriPaquet2023}
P. Clairambault, G. Olimpieri, and F. Paquet.
From thin concurrent games to generalized species of structures.
arXiv:2302.01082, 2023.

\bibitem{VanDenHeuvelPerez2024}
B. van den Heuvel and J. A. P\'erez.
Comparing session type systems derived from linear logic.
arXiv:2401.14763, 2024.

\bibitem{LSR2017}
D. R. Licata, M. Shulman, and M. Riley.
A fibrational framework for substructural and modal logics.
In \emph{2nd International Conference on Formal Structures for Computation and
Deduction (FSCD 2017)}, LIPIcs 84, pp. 25:1--25:22. Schloss Dagstuhl, 2017.
doi:10.4230/LIPIcs.FSCD.2017.25.

\bibitem{RileyFinsterLicata2021}
M. Riley, E. Finster, and D. R. Licata.
Synthetic spectra via a monadic and comonadic modality.
arXiv:2102.04099, 2021.

\end{thebibliography}

\end{document}
